\documentclass{article}

% Recommended, but optional, packages for figures and better typesetting:
\usepackage{microtype}
\usepackage{graphicx}
\usepackage{subcaption}
\usepackage{booktabs} % for professional tables

% hyperref makes hyperlinks in the resulting PDF.
% If your build breaks (sometimes temporarily if a hyperlink spans a page)
% please comment out the following usepackage line and replace
% \usepackage{icml2026} with \usepackage[nohyperref]{icml2026} above.
\usepackage{hyperref}

% Attempt to make hyperref and algorithmic work together better:
\newcommand{\theHalgorithm}{\arabic{algorithm}}

% Use the following line for the initial blind version submitted for review:
%\usepackage{icml2026}

% If accepted, instead use the following line for the camera-ready submission:
\usepackage[accepted]{icml2026}

% Override the ICML "Proceedings of..." footer line: this is a course project,
% not an actual ICML submission. The [accepted] option above is used solely to
% reveal author names and apply camera-ready styling.
\makeatletter
\renewcommand{\ICML@appearing}{\textit{Course project report, CSC~4444G,
Spring~2026, Louisiana State University.}}
\makeatother

% For theorems and such
\usepackage{amsmath}
\usepackage{amssymb}
\usepackage{mathtools}
\usepackage{amsthm}

% if you use cleveref..
\usepackage[capitalize,noabbrev]{cleveref}

%%%%%%%%%%%%%%%%%%%% Theorems-like Environments %%%%%%%%%%%%%%%%%%%%%%%%%%%%%%
% (only used here for completeness; not actually used in our paper)
\theoremstyle{plain}
\newtheorem{theorem}{Theorem}[section]
\newtheorem{proposition}[theorem]{Proposition}
\newtheorem{lemma}[theorem]{Lemma}
\newtheorem{corollary}[theorem]{Corollary}
\theoremstyle{definition}
\newtheorem{definition}[theorem]{Definition}
\newtheorem{assumption}[theorem]{Assumption}
\theoremstyle{remark}
\newtheorem{remark}[theorem]{Remark}
%%%%%%%%%%%%%%%%%%%%%%%%%%%%%%%%%%%%%%%%%%%%%%%%%%%%%%%%%%%%%%%%%%%%%%%%%%%%%%

% The \icmltitle you define below is probably too long as a header.
% Therefore, a short form for the running title is supplied here:
\icmltitlerunning{Guardian Agent: A Bayesian Personalized Weather Hazard Safety Agent}

\begin{document}

\twocolumn[
\icmltitle{Guardian Agent: A Bayesian Personalized\\
           Weather Hazard Safety Agent}

\begin{icmlauthorlist}
\icmlauthor{Erfan Zarafshan}{lsu}
\icmlauthor{Maby Gavilan Abanto}{lsu}
\end{icmlauthorlist}

\icmlaffiliation{lsu}{Division of Computer Science and Engineering,
                       Louisiana State University, Baton Rouge, LA, USA}

\icmlcorrespondingauthor{Erfan Zarafshan}{ezaraf1@lsu.edu}

\icmlkeywords{Bayesian networks, intelligent agents, weather hazards,
              gradient boosting, decision support}

\vskip 0.3in
]

\printAffiliationsAndNotice{}

\begin{abstract}
Existing weather emergency alert systems are reactive and broadcast-style:
the same warning text is sent to every recipient in an affected region,
regardless of their personal vulnerability. We present \emph{Guardian
Agent}, a context-aware AI system that combines real-time weather data,
a gradient-boosted hazard classifier trained on the NOAA Storm Events
Database, and a discrete Bayesian network with ten nodes to produce a
personalized risk posterior, which a deterministic rule planner then
maps to an ordered list of recommended actions. The system was implemented
in Python, deployed as an interactive web demonstration, and evaluated on
hold-out data and four hazard scenarios. The trained classifier reaches
ROC-AUC $= 0.759$ on $21{,}281$ held-out county-day cells, and the
end-to-end system produces visibly distinct recommendations for users
with different vulnerability profiles facing identical hazards,
demonstrating the value of personalized probabilistic reasoning over
generic broadcast warnings.
\end{abstract}

\section{Introduction}
\label{sec:intro}

Severe weather events kill hundreds and cause billions of dollars in
property damage in the United States each year. The dominant
public-warning mechanism, Wireless Emergency Alerts (WEA), pushes the
same short message to every cell phone in a geographic
polygon~\cite{fcc_wea}. Two recipients in the same neighborhood --- one
on the third floor of a sturdy building with a high-clearance vehicle,
the other on the ground floor with a mobility limitation and no car ---
receive identical instructions, even though the appropriate response is
materially different. The result is a well-documented problem of
\emph{alert fatigue} in which recipients learn to discount warnings that
do not appear to apply to them.

This paper addresses that gap by treating the safety-recommendation
problem as a probabilistic inference task over user-specific evidence.
Our system, \textbf{Guardian Agent}, observes the world through public
weather APIs, fuses those observations with a learned threat signal and
a structured user profile, and computes a posterior over four risk
levels \{Low, Moderate, High, Critical\} for that specific user. A rule
planner then converts the argmax of the posterior into an ordered list
of actions chosen from a fixed catalog (notify, shelter, evacuate,
activate flood lights, etc.), each annotated with a human-readable
rationale.

\textbf{Contributions.} (1) An end-to-end goal-based agent
architecture, in the sense of Russell and Norvig
(\citeyear{russell2020ai}, Ch.~2), that integrates an ML classifier and a
hand-engineered Bayesian network into a single perceive-reason-act
loop. (2) A discrete Bayesian network whose conditional probability
tables are specified by domain-knowledge scoring functions rather than
enumerated by hand, scaling to ten nodes while remaining auditable. (3)
A publicly deployed interactive demonstration at
\url{https://guardianagent.streamlit.app/} that lets visitors explore
both real and hypothetical hazards under multiple personal profiles.

\section{Methodology}
\label{sec:methodology}

\subsection{Architecture overview}

Figure~\ref{fig:architecture} shows Guardian Agent's perceive--reason--act
cycle. Each component is a Python module under
\texttt{src/guardian/}; modules are loosely coupled through typed
\texttt{pydantic} dataclasses so that each layer can be tested and
swapped independently.

\begin{figure}[t]
\centering
\fbox{\parbox{0.95\columnwidth}{\centering\small
\textbf{PERCEIVE}\\
NWS API~\cite{nws_api} $+$ OpenWeatherMap~\cite{owm_api}\\
$\downarrow$\\
\textbf{REASON}\\
ML threat classifier (sklearn) $+$\\
Bayesian network (pgmpy) $+$ User profile\\
$\downarrow$\\
\textbf{ACT}\\
Rule planner $\rightarrow$ Console / SMS / Smart-home dispatchers
}}
\caption{Architecture. Each layer is a separate Python module
communicating through typed dataclasses.}
\label{fig:architecture}
\end{figure}

\subsection{Perception layer}

The perception layer queries two complementary sources. The U.S.
National Weather Service API~\cite{nws_api} provides authoritative
hazard \emph{alerts} (event name, severity, urgency, certainty)
indexed by zone code; OpenWeatherMap~\cite{owm_api} provides current
\emph{conditions} (temperature, wind speed, precipitation rate,
humidity). An aggregator merges both into a single
\texttt{WeatherObservation} record. If either source fails, the system
falls back to the other and logs a \texttt{partial observation} flag
rather than crashing.

\subsection{Threat classifier}
\label{sec:classifier}

The classifier is trained to predict, given a (county, day) cell, the
probability that any severe weather event will occur in the
following 24 hours. Training data is the NOAA Storm Events
Database~\cite{noaa_storm_events}: we used 2018--2024 events from
five Gulf Coast states (LA, TX, MS, AL, FL), yielding $106{,}401$
event records. We synthesized one positive cell per event and
sampled negative cells uniformly from the same county-time
distribution, producing $106{,}401$ training rows after class-balancing
weights are applied.

Features are intentionally shallow: month and day-of-year (cyclically
encoded as $\sin/\cos$), county FIPS code (categorical), state
indicator, and a 30-day rolling event count for the same county. We
deliberately avoid radar or atmospheric soundings to keep the
prototype reproducible from a free public dataset; a production system
would naturally consume those.

The model is a Gradient Boosting
Classifier~\cite{friedman2001greedy} from
\texttt{scikit-learn}~\cite{pedregosa2011scikit}. Hyperparameters were
selected by 5-fold stratified cross-validation over a small grid
(\texttt{n\_estimators} $\in \{50,100,200\}$, \texttt{max\_depth} $\in
\{3,5\}$, \texttt{learning\_rate} $\in \{0.05,0.1\}$). Class imbalance
($\sim 1\%$ positive rate) is handled by inverse-frequency
\texttt{sample\_weight}.

\subsection{Bayesian risk engine}
\label{sec:bayes}

The reasoning core is a discrete Bayesian network with 10 nodes
arranged as a polytree, implemented in
\texttt{pgmpy}~\cite{ankan2024pgmpy,ankan2015pgmpy}. The network
fuses three classes of evidence:

\begin{itemize}\itemsep -2pt
\item \textbf{Hazard evidence} from NWS alerts: \texttt{HazardSeverity}
  (5 states), \texttt{Urgency} (5 states).
\item \textbf{Environmental evidence} from OpenWeatherMap:
  \texttt{WindCategory}, \texttt{PrecipCategory} (4 states each).
\item \textbf{Threat-prediction evidence} from the ML classifier:
  \texttt{EmergingThreat} (3 states, derived by binning the classifier
  probability).
\item \textbf{Personal evidence} from the user profile:
  \texttt{HomeFloor}, \texttt{VehicleClearance}, \texttt{MobilityLimited}.
\end{itemize}

These eight nodes feed a latent \texttt{HazardImpact} node, which in
turn determines the query node \texttt{RiskLevel} $\in$
\{Low, Moderate, High, Critical\}. Inference is exact via
Variable Elimination (\citealt{koller2009pgm}, Ch.~9), and completes
in a few milliseconds per query at this network size.

A practical issue with hand-built networks is the combinatorial blowup
of CPT entries (the parents of \texttt{HazardImpact} alone produce
$5\!\times\!5\!\times\!4\!\times\!4\!\times\!3 = 1200$ rows). To avoid
hand-enumeration, we specify each CPT through a small scoring function
that maps parent assignments to a numeric impact score, then bucketizes
that score into the child's discrete states. This collapses CPT
authorship into roughly fifteen lines of readable Python per node and
produces tables that are internally consistent across parent
combinations.

\subsection{Action planner}

The planner is intentionally deterministic. It takes the argmax of the
risk posterior, the user profile, and the active alert and selects
actions from an 8-element catalog
$\{$\textsc{notify\_user}, \textsc{notify\_contacts},
\textsc{recommend\_shelter}, \textsc{recommend\_evacuate},
\textsc{activate\_flood\_lights}, \textsc{sound\_alarm},
\textsc{adjust\_thermostat}, \textsc{log\_only}$\}$ across three
dispatch channels (console, SMS, smart-home). Hazard-specific message
templates are applied per alert event (e.g., a flash-flood template
warns about flooded roadways; a tornado template specifies interior
sheltering). User-preference fields gate non-essential channels:
quiet-hours suppress \emph{moderate}-level notifications but are
overridden for \emph{critical} life-safety actions.

Every action carries a \texttt{rationale} string that is written to the
log alongside the action. This converts the system's behavior into a
post-hoc auditable transcript --- in our judgment a more important
property than maximizing recommendation accuracy on a benchmark, given
the consequences of incorrect emergency advice.

\subsection{Optional LLM explainer}

The deployed web demo offers an optional LLM-powered explainer
(Anthropic Claude or OpenAI GPT, user-supplied API key) that
rephrases the planner's output into conversational guidance. The LLM
\emph{cannot} alter recommendations: it receives the planner's actions
as input and is instructed via a system prompt to neither contradict
nor add to them. This keeps the safety-critical reasoning fully
grounded in the deterministic Bayesian-plus-planner pipeline while
improving the readability of the explanation for end users.

\section{Experimental Results}
\label{sec:results}

\subsection{Operational environment}

The system was developed in Python 3.13 on macOS 14.5
(Apple Silicon). Key libraries:
\texttt{pgmpy} 1.1, \texttt{scikit-learn} 1.6, \texttt{pandas} 2.2,
\texttt{pydantic} 2.10, \texttt{streamlit} 1.32, \texttt{plotly} 5.18.
SMS dispatch is integrated via Twilio in dry-run mode by default; the
deployed web demo enforces dry-run unconditionally to prevent abuse.
External data: NOAA Storm Events 2018--2024 (Gulf Coast states),
NWS-API (live), OpenWeatherMap free-tier (live). Source code at
\url{https://github.com/ErfanZarafshan/guardian-agent}.

\subsection{Classifier evaluation}
\label{sec:classifier-eval}

We evaluate the trained Gradient Boosting Classifier on a stratified
$80/20$ split of the prepared cells.

\begin{table}[h]
\caption{Threat classifier performance on the held-out test split.}
\label{tab:clf}
\centering\small
\begin{tabular}{lr}
\toprule
\textbf{Metric} & \textbf{Value} \\
\midrule
Training cells       & $85{,}120$ \\
Test cells           & $21{,}281$ \\
Positive base rate   & $\approx 0.9\%$ \\
Test ROC-AUC         & $\mathbf{0.759}$ \\
Test accuracy        & $0.696$ \\
Test F\textsubscript{1} & $0.041$ \\
\bottomrule
\end{tabular}
\end{table}

ROC-AUC of $0.759$ is well above chance ($0.500$) but well below the
$0.85\!-\!0.95$ that mature meteorological nowcasting systems achieve
with radar input. We view this as the honest ceiling of a
calendar-and-location-only feature set; the Bayesian fusion downstream
of the classifier compensates partially by leaning more heavily on
NWS alert evidence when the classifier signal is weak.

The very low F\textsubscript{1} reflects the rare-event nature of the
task: the classifier rarely predicts a positive at the default $0.5$
threshold. Operationally we do not threshold; we bin the classifier
probability into three buckets (Low, Moderate, High) and feed those as
discrete \texttt{EmergingThreat} evidence into the Bayesian network,
which is what actually drives downstream behavior.

\subsection{End-to-end personalization}
\label{sec:personalization}

The decisive test of personalization is whether the same hazard
produces different recommendations for different users.
Table~\ref{tab:scenarios} reports posteriors for two reference
profiles --- \emph{Alex} (third-floor apartment, owns vehicle, no
medical limitations) and \emph{Sam} (ground-floor apartment in flood
zone AE, no vehicle, mobility-limited) --- across four scenarios.

\begin{table}[h]
\caption{Posterior probability of each \texttt{RiskLevel} state by
scenario, for two contrasting user profiles. Argmax in \textbf{bold}.
``Argmax flip'' indicates that the planner selects a strictly different
action set for the two users.}
\label{tab:scenarios}
\centering\footnotesize
\begin{tabular}{lcccc}
\toprule
& Low & Mod. & High & Crit. \\
\midrule
\multicolumn{5}{l}{\textit{Clear day}} \\
~~Alex                & \textbf{.81} & .13 & .05 & .01 \\
~~Sam                 & \textbf{.78} & .12 & .07 & .03 \\
\midrule
\multicolumn{5}{l}{\textit{Flash Flood Watch} \quad (argmax flip)} \\
~~Alex                & .26 & \textbf{.37} & .27 & .11 \\
~~Sam                 & .13 & .22 & \textbf{.38} & .28 \\
\midrule
\multicolumn{5}{l}{\textit{Tornado Warning}} \\
~~Alex                & .02 & .12 & .37 & \textbf{.48} \\
~~Sam                 & .00 & .07 & .32 & \textbf{.61} \\
\midrule
\multicolumn{5}{l}{\textit{Tropical Storm Warning} \quad (argmax flip)} \\
~~Alex                & .06 & .19 & \textbf{.39} & .37 \\
~~Sam                 & .02 & .09 & .34 & \textbf{.56} \\
\bottomrule
\end{tabular}
\end{table}

In two of four scenarios (Flash Flood Watch and Tropical Storm
Warning) the posterior argmax differs between Alex and Sam. Because
the planner conditions on the argmax, this difference cascades into a
strictly different action set: Sam receives a shelter recommendation
where Alex receives only a notification, and Sam's higher-priority
contact list is paged where Alex's is not. In the Tornado scenario
both users are placed at Critical, but Sam's posterior assigns
$\sim 13$ percentage points more mass to Critical, reflecting the
amplifying effect of mobility limitation on tornado outcomes.

\subsection{Live agent and dispatch behaviour}

Running the agent against live NWS and OpenWeatherMap data for our
profile address in Baton Rouge, LA on a calm afternoon produced a
posterior of \{Low: $0.485$, Moderate: $0.319$, High: $0.138$,
Critical: $0.057$\}. A subsequent three-cycle continuous run with a
60-second polling interval verified the deduplication logic: the
fingerprint (active-alert events $\times$ argmax) matched between
cycles 2 and 3, suppressing $N \!>\! 0$ non-console actions per
duplicate cycle while keeping the console heartbeat alive. This
behavior is essential for production deployment: without it, an
hour-long flash flood warning would generate twelve identical SMS
messages to each contact.

\subsection{Test coverage and code quality}

The code base ships with $239$ unit and integration tests in $9$ test
files exercising every layer (profile validation $27$ tests, weather
parsers $41$, classifier $33$, Bayesian engine $21$, risk engine
integration $14$, planner $37$, dispatchers $19$, agent loop $15$,
plus smoke tests). All pass on Python 3.13, including end-to-end
tests that mock the weather aggregator and exercise the full
perceive--reason--act loop.

\section{Conclusions and Future Work}
\label{sec:conclusion}

Guardian Agent shows that combining classical probabilistic reasoning
with modern ML and modern web tooling produces a system that
\emph{differentiates} its safety guidance based on user-specific
factors --- the gap that current broadcast emergency alerts leave
open. The Bayesian network is small enough to be auditable and large
enough to capture the interaction between hazard severity, ML threat
prediction, and personal vulnerability; the rule planner is
deterministic and explicable, with every action carrying a written
rationale; and an optional LLM explainer makes the resulting
guidance conversational without subverting the safety-critical
reasoning underneath.

The most valuable extensions, in priority order, would be
(1) replacing the calendar-and-location-only classifier features with
NEXRAD radar reflectivity nowcasts to lift ROC-AUC into the
$0.85\!+$ band; (2) shipping the smart-home dispatcher as a real
Alexa Skills Kit integration so the existing
\texttt{SmartHomeDispatcher} interface can be exercised end-to-end on
a real device; (3) replacing the unconditional dry-run constraint on
the deployed web demo with per-user verified-recipient SMS, allowing
users to receive real notifications while still preventing abuse; and
(4) learning the Bayesian network's CPT scoring functions from
historical NWS alert + outcome data instead of specifying them by
hand. None of these change the architecture --- they replace specific
components behind unchanged interfaces, which is precisely the
property the typed-dataclass module boundaries were designed to
enable.

\bibliography{refs}
\bibliographystyle{icml2026}

\end{document}
